\begin{frame}{Цель и задачи занятия}
  \vspace{0.3em}
  \begin{columns}[T]
    \begin{column}{0.48\textwidth}
      \begin{block}{Цель}
        Запустить индивидуальный проект
        как инженерный артефакт:
        репозиторий, воспроизводимая сборка
        и осмысленный черновик Proposal.
      \end{block}
    \end{column}
    \begin{column}{0.48\textwidth}
      \begin{block}{Задачи}
        \begin{itemize}
          \item разобрать этапы жизненного цикла ПО
                и что из них делаем сегодня;
          \item собрать toolchain: JDK, IDEA, Gradle, Git;
          \item зафиксировать правила работы с AI;
          \item показать три пути и семестровый минимум;
          \item выбрать тему и сформулировать проблему.
        \end{itemize}
      \end{block}
    \end{column}
  \end{columns}
\end{frame}
